% ================================================================

% ================================================================

\section{DLinear}
\label{sec:dlinear}

\subsection{Overview}

DLinear \citep{zeng2023dlinear} is a decomposition-based linear model that
challenges the prevailing trend of applying complex Transformer architectures
to time series forecasting. The core insight is that a time series can be
separated into a slow-moving trend and a fast-moving seasonal component, each
forecast independently with a single linear layer, then recombined. Despite
extreme simplicity --- only 5,892 parameters for the shared-weight variant ---
DLinear matches or exceeds state-of-the-art Transformers on standard benchmarks
while training orders of magnitude faster.

In this project, DLinear is applied as a \textbf{multivariate multi-step}
forecaster: a single model consumes the full $120 \times 11$ input window
(7 weather channels + 4 cyclical time encodings) and outputs a $24 \times 7$
forecast block covering all seven protocol targets simultaneously. A final linear
projection layer maps the 11 input channels to the 7 output targets, allowing the
model to exploit cross-variable dependencies learned from data.

Two variants are evaluated:
\begin{itemize}
  \item \textbf{DLinear-S} (Shared): one pair of trend/seasonal linear layers
        shared across all 11 input channels --- 5,892 parameters.
  \item \textbf{DLinear-I} (Individual): per-channel trend/seasonal linear
        layers --- 63,972 parameters.
\end{itemize}

% ---------------------------------------------------------------
\subsection{Mathematical Formulation}
% ---------------------------------------------------------------

Let $\mathbf{X} \in \mathbb{R}^{L \times C_{\text{in}}}$ be an input window of
$L = 120$ hours and $C_{\text{in}} = 11$ features (7 weather variables + 4
cyclical time encodings). The model produces
$\mathbf{\hat{Y}} \in \mathbb{R}^{H \times C_{\text{out}}}$ with $H = 24$ hours
and $C_{\text{out}} = 7$ target variables.

\subsubsection{Series Decomposition}

A 1D average pooling with kernel size $k = 25$ and stride 1 extracts the trend
while preserving sequence length. For a sequence of length $L$, the first and
last time steps are replicated $p = (k-1)/2 = 12$ times and concatenated to both
ends before pooling --- a technique called \textbf{replication padding} used by
the official implementation \citep{zeng2023dlinear}:
%
\begin{equation}
\mathbf{X}_{\text{padded}} = [\,\underbrace{\mathbf{X}_0, \ldots, \mathbf{X}_0}_{p \text{ times}},\;
        \mathbf{X}_0, \ldots, \mathbf{X}_{L-1},\;
        \underbrace{\mathbf{X}_{L-1}, \ldots, \mathbf{X}_{L-1}}_{p \text{ times}}\,]
        \quad \in \mathbb{R}^{(L + 2p) \times C_{\text{in}}}
\end{equation}
%
\begin{equation}
\mathbf{T} = \text{AvgPool1d}\big(\mathbf{X}_{\text{padded}},\;
            \text{kernel}=k,\; \text{padding}=0\big)
            \quad \in \mathbb{R}^{L \times C_{\text{in}}}
\end{equation}
%
\begin{equation}
\mathbf{S} = \mathbf{X} - \mathbf{T} \quad \in \mathbb{R}^{L \times C_{\text{in}}}
\end{equation}
%
At each position $t$, the trend is the unweighted mean of the $k$ values in the
local window:
%
\begin{equation}
\mathbf{T}_t = \frac{1}{k} \sum_{j=t-p}^{t+p} \mathbf{X}_{\text{padded}, j}
\end{equation}

The kernel size $k=25$ spans just over one full diurnal cycle (24 h), causing
the day--night oscillation to approximately cancel when averaged, leaving only
the multi-day trend. The seasonal component $\mathbf{S}$ captures all remaining
variation: diurnal cycles, short-term fluctuations, and noise.

\paragraph{Why replication padding.}
At the boundaries ($t < p$ and $t \ge L-p$), the kernel extends beyond available
data. Repeating the boundary values --- rather than padding with zeros ---
encodes a \textbf{persistence} assumption: if nothing is known beyond the edge,
the most conservative guess is that conditions remain unchanged. Concretely, at
the leftmost position:
%
\begin{equation}
\mathbf{T}_0 = \frac{\overbrace{\mathbf{X}_0 + \cdots + \mathbf{X}_0}^{p+1 \text{ times}}
        + \mathbf{X}_1 + \cdots + \mathbf{X}_p}{k}
        = \frac{(p+1)\,\mathbf{X}_0 + \sum_{j=1}^{p} \mathbf{X}_j}{k}
\end{equation}
%
All $k=25$ positions in the kernel contain real observations or faithful copies;
no position is underweighted. The first and last 12 positions (24 of 120 h, 20\%
of the lookback window) are affected; the rightmost 12 are the hours closest to
the forecast origin, where the signal is most informative for short-horizon
prediction.

\subsubsection{Dual Linear Projection}

Each component is independently mapped from the 120 h lookback to the 24 h
forecast horizon via a single linear layer:
%
\begin{align}
\mathbf{\hat{Y}}_t &= \mathbf{W}_t \, \mathbf{T} + \mathbf{b}_t
        \quad \in \mathbb{R}^{H \times C_{\text{in}}}
        \qquad \text{(Trend forecast)} \\
\mathbf{\hat{Y}}_s &= \mathbf{W}_s \, \mathbf{S} + \mathbf{b}_s
        \quad \in \mathbb{R}^{H \times C_{\text{in}}}
        \qquad \text{(Seasonal forecast)}
\end{align}
%
where $\mathbf{W}_t, \mathbf{W}_s \in \mathbb{R}^{H \times L}$ and
$\mathbf{b}_t, \mathbf{b}_s \in \mathbb{R}^{H}$.

\begin{itemize}
  \item \textbf{DLinear-S}: a single pair $(\mathbf{W}_t, \mathbf{W}_s)$ is
        shared across all $C_{\text{in}}$ channels.
  \item \textbf{DLinear-I}: each channel $c \in \{1, \ldots, C_{\text{in}}\}$
        has its own pair $\mathbf{W}_t^{(c)}, \mathbf{W}_s^{(c)} \in
        \mathbb{R}^{H \times L}$, implemented as $C_{\text{in}}$ separate
        \texttt{nn.Linear(L, H)} modules.
\end{itemize}

\subsubsection{Combine and Project}
%
\begin{align}
\mathbf{\hat{Y}}_{\text{raw}} &= \mathbf{\hat{Y}}_t + \mathbf{\hat{Y}}_s
        \quad \in \mathbb{R}^{H \times C_{\text{in}}} \\
\mathbf{\hat{Y}} &= \mathbf{\hat{Y}}_{\text{raw}} \, \mathbf{W}_{\text{proj}}
        \quad \in \mathbb{R}^{H \times C_{\text{out}}}
\end{align}
%
$\mathbf{W}_{\text{proj}} \in \mathbb{R}^{C_{\text{in}} \times C_{\text{out}}}$
maps the $C_{\text{in}}$ input channels to the $C_{\text{out}}$ protocol targets.
This projection is linear and data-driven: the model learns weighted combinations
of input channels for each target variable (e.g.\ $wx$ and $wy$ may contribute to
predicting each other, temperature helps predict dew point).

\textbf{Parameter count:}
\begin{itemize}
  \item DLinear-S: $2 \times (LH + H) \times 1 + C_{\text{in}} C_{\text{out}}
        = \mathbf{5{,}892}$
  \item DLinear-I: $2 \times (LH + H) \times C_{\text{in}} + C_{\text{in}}
        C_{\text{out}} = \mathbf{63{,}972}$
\end{itemize}

\subsubsection{Why Decomposition Helps}

The decomposition acts as an \textbf{inductive bias}. The trend branch only needs
to learn a smooth continuation (``is it getting warmer or cooler over the next 24
hours?''), while the seasonal branch only needs to capture the repeating diurnal
pattern (``at what hour does the daily peak occur?''). By separating these two
structurally different tasks, each linear layer solves a simpler problem than one
layer would face on the raw mixed signal.

\textbf{Concrete example (first 48 h of 2021, temperature):}
\begin{table}[H]
\centering
\caption{Decomposition of the first 48 hours of temperature (2021).}
\begin{tabular}{lll}
\toprule
Component & Value range & Meaning \\
\midrule
Raw temp  & $9.4 \to 18.1$\,°C & Day--night oscillation on warming trend \\
Trend ($k=25$) & $5.4 \to 11.4$\,°C & Smooth multi-day warming of $\sim$6\,°C \\
Seasonal & $+4.0$ at noon, $-3.0$ at midnight & Clean 24 h cycle, $\sim\pm 5$\,°C amplitude \\
\bottomrule
\end{tabular}
\end{table}

% ---------------------------------------------------------------
\subsection{Pipeline Overview}
% ---------------------------------------------------------------

Figure~\ref{fig:dlinear-arch} illustrates the DLinear architecture. The model
takes a multivariate input window, decomposes it into trend and seasonal
components via a moving-average kernel, maps each component independently through
linear layers to the forecast horizon, and recombines them with a learned
cross-channel projection.

\begin{figure}[H]
\centering
\fbox{%
\begin{minipage}{0.95\textwidth}
\vspace{6pt}
\[
\underbrace{\mathbf{X} \in \mathbb{R}^{L \times C_{\text{in}}}}_{120\,\text{h} \times 11}
\;\;
\begin{cases}
\;\xrightarrow{\text{AvgPool1d}(k=25)}\;
\underbrace{\mathbf{T} \in \mathbb{R}^{L \times C_{\text{in}}}}_{\text{Trend}}
\;\xrightarrow{\;\mathbf{W}_t \in \mathbb{R}^{H \times L}\;}\;
\underbrace{\mathbf{\hat{Y}}_t}_{\text{Trend forecast}}
&\hspace{-4pt}\searrow \\
&\hspace{-4pt}\Bigg\downarrow \scriptstyle + \;\;
\underbrace{\mathbf{\hat{Y}}}_{\mathbb{R}^{H \times C_{\text{out}}}}
\;\xleftarrow{\;\mathbf{W}_{\text{proj}}\;}
\;\underbrace{\mathbf{\hat{Y}}_{\text{raw}}}_{\mathbb{R}^{H \times C_{\text{in}}}}
\\
\;\xrightarrow{\mathbf{X} - \mathbf{T}}\;
\underbrace{\mathbf{S} \in \mathbb{R}^{L \times C_{\text{in}}}}_{\text{Seasonal}}
\;\xrightarrow{\;\mathbf{W}_s \in \mathbb{R}^{H \times L}\;}\;
\underbrace{\mathbf{\hat{Y}}_s}_{\text{Seasonal forecast}}
&\hspace{-4pt}\nearrow
\end{cases}
\]
\vspace{4pt}
\small
\centering
\textbf{DLinear-S}: one pair $(\mathbf{W}_t,\mathbf{W}_s)$ shared across all $C_{\text{in}}$ channels
\qquad
\textbf{DLinear-I}: per-channel $\mathbf{W}_t^{(c)},\mathbf{W}_s^{(c)}$ for $c=1,\ldots,C_{\text{in}}$
\vspace{4pt}
\end{minipage}}
\caption{DLinear architecture: decomposition $\to$ dual linear projection $\to$ combine $\to$ project.}
\label{fig:dlinear-arch}
\end{figure}

% ---------------------------------------------------------------
\subsection{Training Procedure}
% ---------------------------------------------------------------

\textbf{One model, all targets.} Unlike XGBoost (which trains one regressor per
target), a single DLinear model predicts all seven targets jointly. The final
projection layer $\mathbf{W}_{\text{proj}}$ learns cross-target relationships
directly from the training data.

\textbf{Data usage.} The model trains on the full 2021--2023 sliding windows at
stride 1 (26,137 windows). Every window conditions strictly on observations
before the target horizon, preventing leakage. No subsampling or stride reduction
is applied --- the model is fast enough to use all available windows.

\textbf{Validation and early stopping.} Training minimizes MSE with the Adam
optimizer \citep{kingma2015adam}. A ReduceLROnPlateau scheduler halves the
learning rate when validation loss fails to improve for 5 epochs. Early stopping
terminates training when validation loss does not improve for 10 consecutive
epochs, capped at 50 epochs total. All hyperparameters
(Table~\ref{tab:dlinear-hparams}) are fixed a priori --- there is no grid search
--- and the 2025 test set is never seen during training.

\textbf{Convergence.} On an NVIDIA GeForce RTX 4060 GPU, DLinear-S converges in
$\sim$85 s per seed to a mean validation MSE of $0.3538 \pm 1.6\!\times\!10^{-4}$
(3 seeds). DLinear-I converges in $\sim$233 s per seed to
$0.3451 \pm 5.7\!\times\!10^{-4}$. Both plateau smoothly without divergence:
validation loss slopes over the last 10 epochs are $-1.5 \times 10^{-4}$ (S) and
$-3.6 \times 10^{-4}$ (I) --- both effectively flat. The train--validation loss
gap is small (S: $-0.030$; I: $-0.036$), indicating no overfitting.

\textbf{Multi-seed training.} Per protocol \S8, three seeds
$\{42, 123, 456\}$ are trained and checkpoints saved for each seed per variant.
The purely linear optimization landscape is near-convex, leading to highly
consistent convergence regardless of random initialization: across 3 seeds, MAE
standard deviation is $0.004$ for S and $0.023$ for I.

\begin{table}[H]
\centering
\caption{DLinear training hyperparameters.}
\label{tab:dlinear-hparams}
\begin{tabular}{ll}
\toprule
\textbf{Parameter} & \textbf{Value} \\
\midrule
Optimizer      & Adam \\
Learning rate  & $10^{-3}$ \\
Weight decay   & $10^{-5}$ \\
Loss           & MSE \\
Batch size     & 64 \\
LR scheduler   & ReduceLROnPlateau (factor 0.5, patience 5) \\
Early stopping & Patience 10 epochs, max 50 \\
Kernel size    & 25 \\
Seeds          & $\{42, 123, 456\}$ \\
\bottomrule
\end{tabular}
\end{table}

% ---------------------------------------------------------------
\subsection{Algorithm}
% ---------------------------------------------------------------

\begin{spacing}{1.0}
\begin{algorithm}[H]
\small
\caption{DLinear training and rolling-origin evaluation}
\label{alg:dlinear}
\begin{algorithmic}[1]
\Require Hourly series $S[0..T]$; $L=120$; $H=24$; $V=7$; $k=25$;
         seeds $\{42,123,456\}$
\State \textbf{Preprocess:}
\State \quad Parse time; clip wind $\ge 0$
\State \quad $wx = \text{speed}\cdot\cos(\text{dir})$,\;
        $wy = \text{speed}\cdot\sin(\text{dir})$
\State \quad Add $\sin$/$\cos$ hour-of-day, day-of-year
\State \quad Chronological split: train(2021--23), val(2024), test(2025)
\State \quad Fit StandardScaler on train, transform all
\State \textbf{Generate sliding windows} (stride 1):
\State \quad $X[i] = S[i:i+L] \in \mathbb{R}^{L \times 11}$,\;
        $Y[i] = S[i+L:i+L+H] \in \mathbb{R}^{H \times 7}$
\For{each seed $s \in \{42,123,456\}$}
  \State Init.\ DLinear($k=25$, individual)
  \For{epoch $=1$ to $50$}
    \State $\mathbf{T} = \text{AvgPool1d}(\mathbf{X})$;\;
          $\mathbf{S} = \mathbf{X} - \mathbf{T}$
    \State $\mathbf{\hat{Y}} = (\mathbf{W}_t\mathbf{T} + \mathbf{b}_t)
           + (\mathbf{W}_s\mathbf{S} + \mathbf{b}_s)$
    \State $\mathbf{\hat{Y}} \gets \mathbf{\hat{Y}}\,\mathbf{W}_{\text{proj}}$
    \State $\mathcal{L} = \text{MSE}(\mathbf{\hat{Y}}, \mathbf{Y})$
    \State Adam step; update $\mathbf{W}_t,\mathbf{W}_s,\mathbf{W}_{\text{proj}}$
    \If{val loss stagnant 5 epochs} \State $\eta \gets \eta \times 0.5$ \EndIf
    \If{val loss stagnant 10 epochs} \State \textbf{break} \EndIf
  \EndFor
  \State Save best checkpoint
\EndFor
\State \textbf{Rolling-origin evaluation} (test 2025, step 24 h, $N\approx 360$):
\For{each origin $t$}
  \State $\mathbf{\hat{Y}}_t = \text{model}\big(S[t-L:t]\big)
         \in \mathbb{R}^{24 \times 7}$
  \State Inverse-transform $\mathbf{\hat{Y}}_t$ to real scale
\EndFor
\State \textbf{Aggregate:} overall MAE/RMSE; per-variable; per-horizon;
       bootstrap 95\% CI
\end{algorithmic}
\end{algorithm}
\end{spacing}

% ---------------------------------------------------------------
\subsection{Experimental Results}
% ---------------------------------------------------------------

Evaluated by rolling-origin over the 2025 test year ($N = 360$ non-overlapping
origins, step = 24 h, real scale). Both variants substantially outperform the
naive repeat-last-value baseline (MAE = 6.71). Values are ensemble-mean across
3 seeds $\{42, 123, 456\}$.

\subsubsection{Overall Performance}

\begin{table}[H]
\centering
\caption{Overall test-set performance (2025, real scale, 3-seed ensemble).}
\label{tab:dlinear-overall}
\begin{tabular}{lcccccr}
\toprule
\textbf{Model} & \textbf{MAE} & \textbf{RMSE} & \textbf{MAE 95\% CI} &
\textbf{RMSE 95\% CI} & \textbf{Params} & \textbf{Train} \\
\midrule
DLinear-S & 5.79 & 12.27 & [5.59, 5.99] & [11.73, 12.79] & 5,892   & 85 s \\
DLinear-I & 5.64 & 12.19 & [5.41, 5.82] & [11.63, 12.70] & 63,972  & 233 s \\
\bottomrule
\end{tabular}
\end{table}

DLinear-I improves over DLinear-S by $+2.7\%$ (MAE) and $+0.6\%$ (RMSE), at the
cost of $10.9\times$ more parameters. This aligns with the original paper's
findings \citep{zeng2023dlinear}: individual weights remove the forced compromise
of sharing one temporal pattern across physically diverse channels, allowing each
variable to learn its own dynamics. The RMSE/MAE ratio is $\sim$2.1 for both
variants, indicating occasional large errors driven by cloud cover and wind
extremes. Bootstrap 95\% CIs (1000 resamples) are narrow (width $<4\%$),
confirming stable metric estimates.

\subsubsection{Per-Variable Error}

\begin{table}[H]
\centering
\caption{Per-variable MAE (real scale, 2025).}
\label{tab:dlinear-pervar}
\begin{tabular}{lrrrr}
\toprule
\textbf{Target} & \textbf{S MAE} & \textbf{I MAE} & \textbf{Unit} & \textbf{I vs.\ S} \\
\midrule
temperature       & 1.22 & 1.18 & °C   & $+2.7\%$ \\
dew\_point        & 1.28 & 1.19 & °C   & $+6.4\%$ \\
surface\_pressure & 1.33 & 1.13 & hPa  & $+15.2\%$ \\
relative\_humidity& 5.70 & 5.64 & \%   & $+1.1\%$ \\
wx                & 4.47 & 4.39 & km/h & $+1.8\%$ \\
wy                & 3.10 & 3.07 & km/h & $+1.1\%$ \\
cloud\_cover      & 23.41& 23.00& \%   & $+1.8\%$ \\
\bottomrule
\end{tabular}
\end{table}

DLinear-I improves on every variable. The largest gain is surface pressure
($+15.2\%$), whose slow-varying, weak-diurnal-cycle dynamics differ fundamentally
from temperature and humidity (strong 24 h cycle). Shared weights force a
compromise across these diverse patterns; individual weights remove this
constraint, confirming the paper's channel-independence hypothesis.

Cloud cover remains the hardest target across all models (MAE $\sim$23\%), driven
by its bimodal distribution (clustered at 0\% and 100\%) and weak linear
correlation with other variables.

\subsubsection{Per-Horizon Error Growth}

Error grows approximately 2.0$\times$ from $h=1$ to $h=24$ (DLinear-S):

\begin{table}[H]
\centering
\caption{Per-horizon MAE (DLinear-S, 2025 test).}
\begin{tabular}{lccccccc}
\toprule
Horizon & $h=1$ & $h=4$ & $h=8$ & $h=12$ & $h=16$ & $h=20$ & $h=24$ \\
\midrule
MAE & 3.46 & 4.58 & 5.07 & 5.77 & 6.75 & 6.82 & 7.03 \\
\bottomrule
\end{tabular}
\end{table}

The error growth is approximately linear with horizon. The plateau between
$h=20$ and $h=24$ suggests the model's uncertainty saturates beyond roughly
20 hours --- additional lookback or auxiliary features may be needed for
long-horizon forecasts.

\subsubsection{Diagnostic Summary}

\begin{table}[H]
\centering
\caption{Protocol \S6 diagnostic results (2025 test).}
\small
\begin{tabular}{l p{2.3cm} p{2.3cm} p{5.2cm}}
\toprule
\textbf{Check} & \textbf{DLinear-S} & \textbf{DLinear-I} &
\textbf{Interpretation} \\
\midrule
ADF stationarity
  & PASS ($p\approx 0$)
  & PASS ($p\approx 0$)
  & All 7 targets stationary, $d=0$ \\
Stability
  & PASS (slope $-1.5\!\times\!10^{-4}$)
  & PASS (slope $-3.6\!\times\!10^{-4}$)
  & Flat val loss, fully converged \\
Ljung--Box (lag 24)
  & 76/168 (45\%)
  & 90/168 (54\%)
  & Autocorrelation at short horizons ($h=1..6$); pressure
    \& dew\_point worst per Protocol \S11 \\
Normality
  & FAIL ($p\approx 0$)
  & FAIL ($p\approx 0$)
  & Expected --- cloud bimodal, wind heavy-tailed \\
Forecast int.\ (95\%)
  & Temp $\pm 1.85\!\to\!\pm 2.90$\,°C
  & Temp $\pm 1.64\!\to\!\pm 2.76$\,°C
  & Monotonic widening; pressure widens $\times 5.3$ in I \\
Bootstrap CI (MAE)
  & $\pm 3.5\%$
  & $\pm 3.7\%$
  & Narrow ($<5\%$), reliable metrics \\
\bottomrule
\end{tabular}
\end{table}

\textbf{ADF:} All seven target variables are stationary at the 1\% level --- no
differencing is needed.

\textbf{Ljung--Box:} Residual autocorrelation is present in roughly half of the
168 horizon-variable pairs, concentrated at short horizons ($h=1..6$) and in
pressure and dew\_point --- variables Protocol \S11 identifies as having known
long memory. DLinear-I's individual weights partially mitigate this (54\% pass
vs.\ 45\% for S).

\textbf{Residual normality:} Universally rejected, as expected. Cloud residuals
are U-shaped (bimodal target), wind residuals are heavy-tailed (gust events).
Gaussian-based confidence intervals would be unreliable; empirical quantile-based
intervals are used instead.

\textbf{Forecast intervals:} All variables show monotonic widening from $h=1$ to
$h=24$. Pressure widens most aggressively ($\times 5.3$ in DLinear-I) due to its
long-memory structure --- the model's uncertainty about pressure grows faster
than for temperature ($\times 1.6$).

\textbf{Bootstrap CI:} MAE confidence interval width under $\pm 4\%$ for both
variants confirms stable metric estimates.

% ---------------------------------------------------------------
\subsection{Analysis}
% ---------------------------------------------------------------

\subsubsection{What DLinear Does Well}

\textbf{Variables with strong diurnal cycles} (temperature, dew point, relative
humidity) are forecast tightly. Temperature MAE of 1.18\,°C (DLinear-I) is
competitive with the best models in this study. This is a direct consequence of
the decomposition design: the seasonal branch is purpose-built to capture the
24 h oscillation, and the sinusoidal time encodings give the model explicit
knowledge of where it sits within the cycle.

\textbf{Surface pressure} benefits most from individual weights ($+15.2\%$ I over
S). Unlike temperature, pressure does not follow a clean 24 h sinusoid --- its
variations are driven by synoptic-scale weather systems evolving over multiple
days. Shared weights, dominated by the strong diurnal signals of temperature and
humidity, are suboptimal for pressure. Individual weights free pressure's linear
layers to learn its own, slower temporal dynamics.

\subsubsection{What DLinear Struggles With}

\textbf{Cloud cover} is the hardest variable (MAE $\sim$23\%) across both
variants and all models in this study. Three factors conspire: (a) cloud cover is
fundamentally bimodal (clear = 0\%, overcast = 100\%), making MSE poorly suited;
(b) it has the weakest linear correlation with other variables; (c) its
transitions (clear $\leftrightarrow$ overcast) are often sharp and driven by
unobserved atmospheric processes. A dedicated classification-based approach may
be warranted.

\textbf{Wind components} ($wx$, $wy$) have moderate MAE (3--4 km/h) but high
RMSE/MAE ratios ($\sim$1.4), indicating occasional large misses during gust
events. Wind gusts are inherently stochastic at hourly resolution --- the 120 h
lookback contains limited predictive signal for a specific gust 24 h ahead.

\textbf{Residual autocorrelation} in dew\_point and pressure suggests the 120 h
lookback may be insufficient for variables with long memory. Increasing the
lookback (e.g.\ to 168 h) or adding explicit lag features could capture more of
this structure.

\subsubsection{S vs.\ I Trade-off in Practice}

DLinear-I's $+2.7\%$ MAE improvement comes at $10.9\times$ the parameters.
Whether this is worthwhile depends on context:

\begin{itemize}
  \item \textbf{For model comparison:} DLinear-I should be the reported variant
        --- it represents the best achievable with the DLinear architecture and
        follows the principle of giving each model its fair configuration.
  \item \textbf{For deployment:} DLinear-S is the pragmatic choice --- 85 s
        training, 5,892 parameters, and only marginally worse accuracy. The
        $2.7\%$ gap is within the bootstrap CI width ($\sim 3.5\%$).
  \item \textbf{For interpretability:} DLinear-S is strongly preferred --- its
        single pair of $\mathbf{W}_t, \mathbf{W}_s$ matrices ($24 \times 120$
        each) can be visualized as heatmaps showing what past hours influence
        each future hour, for trend vs.\ seasonal branches separately.
\end{itemize}
